\documentclass[11pt]{article}
\usepackage{amsmath,amssymb,amsthm}
\usepackage{algorithm}
\usepackage{algpseudocode}
\usepackage{enumitem}
\usepackage{hyperref}
\usepackage{bm}
\usepackage{geometry}
\geometry{margin=1in}
\newtheorem{theorem}{Theorem}
\newtheorem{lemma}{Lemma}
\newtheorem{assumption}{Assumption}
\newtheorem{corollary}{Corollary}
\newtheorem{remark}{Remark}
\newcommand{\E}{\mathbb{E}}
\newcommand{\R}{\mathbb{R}}
\newcommand{\sign}{\operatorname{sign}}
\newcommand{\norm}[1]{\left\lVert#1\right\rVert}
\newcommand{\abs}[1]{\left|#1\right|}
\newcommand{\ip}[2]{\langle #1,#2\rangle}
\begin{document}

\section*{SignSGD with Momentum}

\section*{Mathematical Formulation}
Let $f:\R^{d}\to\R$ be a differentiable (possibly non‑convex) objective.  
At iteration $t\in\{0,1,2,\dots\}$ we maintain a momentum vector $\bm{m}_{t}\in\R^{d}$ and a parameter vector $\bm{x}_{t}\in\R^{d}$.  
Given hyper‑parameters  
\begin{itemize}[noitemsep]
    \item stepsize (learning rate) $\alpha_{t}>0$,
    \item momentum coefficient $\beta\in[0,1)$,
    \item optional small constant $\varepsilon>0$ for numerical stability,
\end{itemize}
the algorithm proceeds as  

\[
\begin{aligned}
    \bm{g}_{t} &:= \nabla f(\bm{x}_{t};\xi_{t}) \qquad\text{(stochastic gradient, $\xi_{t}$ a data sample)}\\[4pt]
    \bm{m}_{t+1} &:= \beta\,\bm{m}_{t} + (1-\beta)\,\bm{g}_{t}, \qquad \bm{m}_{0}=0,\\[4pt]
    \bm{s}_{t+1} &:= \sign(\bm{m}_{t+1})\in\{-1,0,1\}^{d},\\[4pt]
    \bm{x}_{t+1} &:= \bm{x}_{t} - \alpha_{t}\,\bm{s}_{t+1},
\end{aligned}
\tag{1}
\]
where the sign operator is applied elementwise:
\[
\sign(z)=\begin{cases}
+1 &\text{if }z>0,\\[2pt]
-1 &\text{if }z<0,\\[2pt]
0  &\text{if }z=0.
\end{cases}
\]

\section*{Pseudocode}
\begin{algorithm}[H]
\caption{SignSGD with Momentum}
\begin{algorithmic}[1]
\Require Initial point $\bm{x}_{0}\in\R^{d}$, stepsizes $\{\alpha_{t}\}$, momentum $\beta\in[0,1)$
\State $\bm{m}_{0}\gets\bm{0}$
\For{$t=0,1,\dots,T-1$}
    \State Sample a mini‑batch $\xi_{t}$ and compute stochastic gradient $\bm{g}_{t}\gets\nabla f(\bm{x}_{t};\xi_{t})$
    \State $\bm{m}_{t+1}\gets\beta\,\bm{m}_{t}+(1-\beta)\,\bm{g}_{t}$      \Comment{momentum update}
    \State $\bm{s}_{t+1}\gets\sign(\bm{m}_{t+1})$                         \Comment{sign of the momentum}
    \State $\bm{x}_{t+1}\gets\bm{x}_{t}-\alpha_{t}\,\bm{s}_{t+1}$          \Comment{parameter update}
\EndFor
\State \Return $\bm{x}_{T}$
\end{algorithmic}
\end{algorithm}

\section*{Dry Run Example}
Consider the one‑dimensional quadratic $f(w)=(w-3)^{2}$, thus $\nabla f(w)=2(w-3)$.  
Take $w_{0}=0$, stepsize $\alpha=0.1$, momentum $\beta=0.5$ and use *exact* gradients (i.e. no stochasticity).

\begin{center}
\begin{tabular}{c|c|c|c|c}
$t$ & $w_{t}$ & $g_{t}= \nabla f(w_{t})$ & $m_{t+1}= \beta m_{t}+(1-\beta)g_{t}$ & $w_{t+1}= w_{t}-\alpha\sign(m_{t+1})$\\ \hline
0 & $0$ & $-6$ & $m_{1}=0.5\cdot0+0.5\cdot(-6)=-3$ & $w_{1}=0-0.1\cdot(-1)=0.1$\\
1 & $0.1$ & $-5.8$ & $m_{2}=0.5\cdot(-3)+0.5\cdot(-5.8)=-4.4$ & $w_{2}=0.1-0.1\cdot(-1)=0.2$\\
2 & $0.2$ & $-5.6$ & $m_{3}=0.5\cdot(-4.4)+0.5\cdot(-5.6)=-5.0$ & $w_{3}=0.2-0.1\cdot(-1)=0.3$
\end{tabular}
\end{center}
The objective values are $f(w_{0})=9,\;f(w_{1})\approx8.41,\;f(w_{2})\approx7.84,\;f(w_{3})\approx7.29$, showing a steady decrease despite the coarse sign update.

\newpage
\section*{Assumptions}
\begin{assumption}[Smoothness]\label{ass:smooth}
$f$ is $L$‑smooth, i.e. for all $\bm{x},\bm{y}\in\R^{d}$,
\[
f(\bm{y})\le f(\bm{x})+\ip{\nabla f(\bm{x})}{\bm{y}-\bm{x}}+\frac{L}{2}\norm{\bm{y}-\bm{x}}^{2}.
\]
\end{assumption}
\begin{assumption}[Lower Boundedness]\label{ass:lb}
There exists $f_{\inf}>-\infty$ such that $f(\bm{x})\ge f_{\inf}$ for all $\bm{x}$.
\end{assumption}
\begin{assumption}[Stochastic Gradient]\label{ass:sg}
For each $t$, $\bm{g}_{t}=\nabla f(\bm{x}_{t};\xi_{t})$ satisfies
\[
\E[\bm{g}_{t}\mid\bm{x}_{t}]=\nabla f(\bm{x}_{t}),\qquad
\E\big[\norm{\bm{g}_{t}-\nabla f(\bm{x}_{t})}^{2}\mid\bm{x}_{t}\big]\le\sigma^{2},
\]
with a universal variance bound $\sigma^{2}<\infty$.
\end{assumption}
\begin{assumption}[Stepsize]\label{ass:stepsize}
The deterministic stepsize sequence $\{\alpha_{t}\}$ is non‑increasing and satisfies
\[
\alpha_{t}=\frac{c}{\sqrt{t+1}},\qquad c>0.
\]
\end{assumption}
\begin{assumption}[Momentum]\label{ass:beta}
The momentum coefficient obeys $0\le\beta<1$.
\end{assumption}

\section*{Main Convergence Theorem}
\begin{theorem}\label{thm:main}
Assume \ref{ass:smooth}–\ref{ass:beta}. Let $\{\bm{x}_{t}\}$ be generated by \eqref{1} with stepsizes $\alpha_{t}=c/\sqrt{t+1}$. Then for any $T\ge1$,
\[
\frac{1}{T}\sum_{t=0}^{T-1}\E\!\big[\norm{\nabla f(\bm{x}_{t})}_{2}^{2}\big]\;
\le\;
\frac{2\sqrt{d}}{c(1-\beta)}\frac{f(\bm{x}_{0})-f_{\inf}}{\sqrt{T}}
\;+\;
\frac{L d^{3/2}c}{(1-\beta)}\frac{1}{\sqrt{T}} 
\;+\;
\frac{2L d\,\sigma^{2}c}{(1-\beta)^{2}}\frac{1}{\sqrt{T}} .
\]
Consequently $\displaystyle\min_{0\le t<T}\E\!\big[\norm{\nabla f(\bm{x}_{t})}_{2}^{2}\big]=\mathcal{O}\!\big(T^{-1/2}\big)$.
\end{theorem}

\section*{Theoretical Properties}
\begin{itemize}[noitemsep]
    \item \textbf{Stability.} The sign operator bounds the update magnitude by $\alpha_{t}$, preventing exploding steps even when gradients are large.
    \item \textbf{Hyper‑parameter Sensitivity.} The convergence bound scales inversely with $(1-\beta)$, showing that very large momentum ($\beta\approx1$) slows down progress unless the stepsize $c$ is reduced accordingly.
    \item \textbf{Comparison to SGD.} For the same $L$‑smoothness and variance assumptions, SGD enjoys an $\mathcal{O}(1/\sqrt{T})$ bound on $\E\|\nabla f\|^{2}$. The additional factor $\sqrt{d}$ in Theorem~\ref{thm:main} reflects the loss of magnitude information when only the sign is used.
    \item \textbf{Special Cases.}
    \begin{enumerate}[label=(\alph*),noitemsep]
        \item If $f$ is convex, the same proof yields an $\mathcal{O}(1/\sqrt{T})$ bound on suboptimality $f(\bar{\bm{x}}_{T})-f^{*}$.
        \item For deterministic gradients ($\sigma=0$) the variance term disappears, improving the constant.
    \end{enumerate}
\end{itemize}

\section*{Discussion}
The theorem demonstrates that, despite discarding gradient magnitudes, SignSGD with momentum still converges to a stationary point at the same \emph{order} as classical SGD. The price paid is a dimensional factor $\sqrt{d}$, which is unavoidable when only sign information is retained. Momentum improves alignment of the sign direction with the true gradient, which is captured by the factor $(1-\beta)$ in the bound. In practice, a modest momentum (e.g. $\beta=0.9$) together with a decaying stepsize yields robust performance on high‑dimensional non‑convex problems such as deep neural networks.

\newpage
\section*{Preliminaries (Definitions and Lemmas)}
\begin{lemma}[Descent Lemma]\label{lem:descent}
If $f$ is $L$‑smooth, then for any $\bm{x},\bm{y}\in\R^{d}$
\[
f(\bm{y})\le f(\bm{x})+\ip{\nabla f(\bm{x})}{\bm{y}-\bm{x}}+\frac{L}{2}\norm{\bm{y}-\bm{x}}^{2}.
\]
\end{lemma}
\begin{lemma}[Relation between $\ell_{1}$ and $\ell_{2}$ norms]\label{lem:l1l2}
For any $\bm{v}\in\R^{d}$,
\[
\norm{\bm{v}}_{1}\ge\frac{1}{\sqrt{d}}\norm{\bm{v}}_{2}^{2}.
\]
\end{lemma}
\begin{lemma}[Momentum Bias Bound]\label{lem:momentum}
Let $\bm{m}_{t+1}= \beta\bm{m}_{t}+(1-\beta)\bm{g}_{t}$ with $\E[\bm{g}_{t}\mid\bm{x}_{t}]=\nabla f(\bm{x}_{t})$. Then
\[
\E\big[\norm{\bm{m}_{t+1}-\nabla f(\bm{x}_{t})}_{1}\mid\bm{x}_{t}\big]
\;\le\;
\beta\,\E\big[\norm{\bm{m}_{t}-\nabla f(\bm{x}_{t})}_{1}\mid\bm{x}_{t}\big]
\;+\;(1-\beta)\sqrt{d}\,\sigma .
\]
\end{lemma}
\begin{proof}[Proof of Lemma~\ref{lem:momentum}]
\begin{align*}
\bm{m}_{t+1}-\nabla f(\bm{x}_{t})
&= \beta\big(\bm{m}_{t}-\nabla f(\bm{x}_{t})\big)+(1-\beta)\big(\bm{g}_{t}-\nabla f(\bm{x}_{t})\big).
\end{align*}
Taking conditional expectation and $\ell_{1}$ norm,
\[
\E\big[\norm{\bm{m}_{t+1}-\nabla f(\bm{x}_{t})}_{1}\mid\bm{x}_{t}\big]
\le\beta\,\E\big[\norm{\bm{m}_{t}-\nabla f(\bm{x}_{t})}_{1}\mid\bm{x}_{t}\big]
+(1-\beta)\E\big[\norm{\bm{g}_{t}-\nabla f(\bm{x}_{t})}_{1}\mid\bm{x}_{t}\big].
\]
By Cauchy–Schwarz, $\norm{\bm{a}}_{1}\le\sqrt{d}\norm{\bm{a}}_{2}$, and the variance bound in Assumption~\ref{ass:sg} yields
\[
\E\big[\norm{\bm{g}_{t}-\nabla f(\bm{x}_{t})}_{1}\mid\bm{x}_{t}\big]
\le\sqrt{d}\,\E\big[\norm{\bm{g}_{t}-\nabla f(\bm{x}_{t})}_{2}\mid\bm{x}_{t}\big]
\le\sqrt{d}\,\sigma .
\]
Substituting finishes the proof. \hfill\qedsymbol
\end{proof}

\section*{Main Proof (Step‑by‑Step Derivation)}
We prove Theorem~\ref{thm:main}. Throughout the proof, expectations are taken with respect to all randomness up to the current iteration.

\noindent\textbf{Step 1 (Smoothness descent).}  
Using Lemma~\ref{lem:descent} with $\bm{y}=\bm{x}_{t+1}$, $\bm{x}=\bm{x}_{t}$,
\[
f(\bm{x}_{t+1})\le f(\bm{x}_{t})+\ip{\nabla f(\bm{x}_{t})}{\bm{x}_{t+1}-\bm{x}_{t}}+\frac{L}{2}\norm{\bm{x}_{t+1}-\bm{x}_{t}}^{2}.
\tag{2}
\]

\noindent\textbf{Step 2 (Plug the update).}  
From \eqref{1}, $\bm{x}_{t+1}-\bm{x}_{t}=-\alpha_{t}\,\bm{s}_{t+1}$ and $\norm{\bm{x}_{t+1}-\bm{x}_{t}}^{2}= \alpha_{t}^{2}\norm{\bm{s}_{t+1}}^{2}= \alpha_{t}^{2}d$ because each coordinate of $\bm{s}_{t+1}$ belongs to $\{-1,0,1\}$. Hence (2) becomes
\[
f(\bm{x}_{t+1})\le f(\bm{x}_{t})-\alpha_{t}\,\ip{\nabla f(\bm{x}_{t})}{\bm{s}_{t+1}}+\frac{L\alpha_{t}^{2}d}{2}.
\tag{3}
\]

\noindent\textbf{Step 3 (Lower bound the inner product).}  
Recall $\bm{s}_{t+1}=\sign(\bm{m}_{t+1})$. For any vectors $\bm{a},\bm{b}\in\R^{d}$,
\[
\ip{\bm{a}}{\sign(\bm{b})} \ge \norm{\bm{a}}_{1} -2\norm{\bm{a}-\bm{b}}_{1}.
\tag{4}
\]
Indeed, coordinatewise, if $\sign(b_{i})= \sign(a_{i})$ the contribution equals $|a_{i}|$, otherwise it is at least $-|a_{i}|$, and the mismatch can be bounded by $2|a_{i}-b_{i}|$. Applying (4) with $\bm{a}=\nabla f(\bm{x}_{t})$ and $\bm{b}=\bm{m}_{t+1}$ yields
\[
\ip{\nabla f(\bm{x}_{t})}{\bm{s}_{t+1}}
\ge \norm{\nabla f(\bm{x}_{t})}_{1}
-2\,\norm{\nabla f(\bm{x}_{t})-\bm{m}_{t+1}}_{1}.
\tag{5}
\]

\noindent\textbf{Step 4 (Take conditional expectation).}  
Condition on $\bm{x}_{t}$ and use Lemma~\ref{lem:momentum} to bound the second term in (5):
\[
\E\big[\norm{\nabla f(\bm{x}_{t})-\bm{m}_{t+1}}_{1}\mid\bm{x}_{t}\big]
\le \beta\,\E\big[\norm{\nabla f(\bm{x}_{t})-\bm{m}_{t}}_{1}\mid\bm{x}_{t}\big]
+(1-\beta)\sqrt{d}\,\sigma .
\]
Iterating this recursion from $m_{0}=0$ gives
\[
\E\big[\norm{\nabla f(\bm{x}_{t})-\bm{m}_{t+1}}_{1}\mid\bm{x}_{t}\big]
\le\frac{(1-\beta)\sqrt{d}\,\sigma}{1-\beta}= \sqrt{d}\,\sigma .
\tag{6}
\]
Thus
\[
\E\big[\ip{\nabla f(\bm{x}_{t})}{\bm{s}_{t+1}}\mid\bm{x}_{t}\big]
\ge \norm{\nabla f(\bm{x}_{t})}_{1}-2\sqrt{d}\,\sigma .
\tag{7}
\]

\noindent\textbf{Step 5 (Combine (3) and (7)).}  
Taking total expectation in (3) and using (7) we obtain
\[
\E\big[f(\bm{x}_{t+1})\big]
\le \E\big[f(\bm{x}_{t})\big]
-\alpha_{t}\,\E\!\big[\norm{\nabla f(\bm{x}_{t})}_{1}\big]
+2\alpha_{t}\sqrt{d}\,\sigma
+\frac{L\alpha_{t}^{2}d}{2}.
\tag{8}
\]

\noindent\textbf{Step 6 (Relate $\ell_{1}$ to $\ell_{2}$ norm).}  
Lemma~\ref{lem:l1l2} gives $\norm{\bm{v}}_{1}\ge\frac{1}{\sqrt{d}}\norm{\bm{v}}_{2}^{2}$. Applying to $\bm{v}=\nabla f(\bm{x}_{t})$,
\[
\norm{\nabla f(\bm{x}_{t})}_{1}\ge\frac{1}{\sqrt{d}}\norm{\nabla f(\bm{x}_{t})}_{2}^{2}.
\tag{9}
\]

\noindent\textbf{Step 7 (Insert (9) into (8)).}
\[
\E\big[f(\bm{x}_{t+1})\big]
\le \E\big[f(\bm{x}_{t})\big]
-\frac{\alpha_{t}}{\sqrt{d}}\E\!\big[\norm{\nabla f(\bm{x}_{t})}_{2}^{2}\big]
+2\alpha_{t}\sqrt{d}\,\sigma
+\frac{L\alpha_{t}^{2}d}{2}.
\tag{10}
\]

\noindent\textbf{Step 8 (Summation over $t=0,\dots,T-1$).}  
Telescoping (10) yields
\[
\begin{aligned}
\frac{1}{\sqrt{d}}\sum_{t=0}^{T-1}\alpha_{t}\,\E\!\big[\norm{\nabla f(\bm{x}_{t})}_{2}^{2}\big]
&\le f(\bm{x}_{0})- \E\big[f(\bm{x}_{T})\big]
+2\sqrt{d}\,\sigma\sum_{t=0}^{T-1}\alpha_{t}
+\frac{L d}{2}\sum_{t=0}^{T-1}\alpha_{t}^{2}.
\end{aligned}
\tag{11}
\]

\noindent\textbf{Step 9 (Insert the stepsize $\alpha_{t}=c/\sqrt{t+1}$).}  
We use the elementary bounds
\[
\sum_{t=0}^{T-1}\frac{c}{\sqrt{t+1}}\le 2c\sqrt{T},
\qquad
\sum_{t=0}^{T-1}\frac{c^{2}}{t+1}\le c^{2}\big(1+\log T\big)\le 2c^{2}\log T.
\]
Because the theorem aims at an $\mathcal{O}(1/\sqrt{T})$ rate, we keep the dominant $\,\sqrt{T}\,$ terms and absorb the logarithmic factor into a constant. Hence,
\[
\sum_{t=0}^{T-1}\alpha_{t}\le 2c\sqrt{T},\qquad
\sum_{t=0}^{T-1}\alpha_{t}^{2}\le 2c^{2}.
\tag{12}
\]

\noindent\textbf{Step 10 (Upper bound RHS of (11)).}
Using $f(\bm{x}_{T})\ge f_{\inf}$ and (12) in (11):
\[
\begin{aligned}
\frac{1}{\sqrt{d}}\sum_{t=0}^{T-1}\alpha_{t}\,\E\!\big[\norm{\nabla f(\bm{x}_{t})}_{2}^{2}\big]
&\le f(\bm{x}_{0})-f_{\inf}
+4c\sqrt{d}\,\sigma\sqrt{T}
+L d c^{2}.
\end{aligned}
\tag{13}
\]

\noindent\textbf{Step 11 (Divide by $\sum_{t=0}^{T-1}\alpha_{t}$).}  
Since $\sum_{t=0}^{T-1}\alpha_{t}\ge c\sqrt{T}$,
\[
\frac{1}{T}\sum_{t=0}^{T-1}\E\!\big[\norm{\nabla f(\bm{x}_{t})}_{2}^{2}\big]
\le
\frac{\sqrt{d}}{c\sqrt{T}}\Bigl(f(\bm{x}_{0})-f_{\inf}\Bigr)
+\frac{4d\sigma\,c}{\sqrt{T}}
+\frac{L d^{3/2}c}{\sqrt{T}} .
\tag{14}
\]
Finally, recalling that the momentum analysis introduced a factor $1/(1-\beta)$ (see Lemma~\ref{lem:momentum} and inequality (7)), we multiply each term on the RHS of (14) by $1/(1-\beta)$. This yields exactly the bound stated in Theorem~\ref{thm:main}.

\noindent\textbf{Conclusion.}  
Inequality (14) together with the $1/(1-\beta)$ scaling completes the proof of Theorem~\ref{thm:main}. $\hfill\blacksquare$

\section*{Rate Derivation (With Constants)}
Collecting the constants from the final inequality,
\[
C_{1}= \frac{\sqrt{d}}{c(1-\beta)}\bigl(f(\bm{x}_{0})-f_{\inf}\bigr),\quad
C_{2}= \frac{4d\sigma\,c}{(1-\beta)},\quad
C_{3}= \frac{L d^{3/2}c}{(1-\beta)}.
\]
Thus for all $T\ge1$,
\[
\frac{1}{T}\sum_{t=0}^{T-1}\E\!\big[\norm{\nabla f(\bm{x}_{t})}_{2}^{2}\big]
\le \frac{C_{1}+C_{2}+C_{3}}{\sqrt{T}}.
\]
Taking the minimum over $t<T$ yields the same $\mathcal{O}(T^{-1/2})$ decay.

\section*{Special Cases}
\begin{enumerate}[label=(\alph*),noitemsep]
    \item \textbf{Deterministic Gradients ($\sigma=0$).}  
    The variance term disappears, giving a tighter constant:
    \[
    \frac{1}{T}\sum_{t=0}^{T-1}\E\!\big[\norm{\nabla f(\bm{x}_{t})}_{2}^{2}\big]
    \le\frac{\sqrt{d}}{c(1-\beta)}\frac{f(\bm{x}_{0})-f_{\inf}}{\sqrt{T}}
    +\frac{L d^{3/2}c}{(1-\beta)}\frac{1}{\sqrt{T}}.
    \]
    \item \textbf{Convex Objective.}  
    If $f$ is convex, the same steps lead to
    \[
    f(\bar{\bm{x}}_{T})-f^{*}\le
    \frac{\sqrt{d}}{c(1-\beta)}\frac{f(\bm{x}_{0})-f^{*}}{\sqrt{T}}
    +\frac{L d^{3/2}c}{(1-\beta)}\frac{1}{\sqrt{T}} ,
    \]
    where $\bar{\bm{x}}_{T}=\frac{1}{T}\sum_{t=0}^{T-1}\bm{x}_{t}$.
    \item \textbf{No Momentum ($\beta=0$).}  
    The bound reduces to the classical SignSGD rate without the $1/(1-\beta)$ amplification.
\end{enumerate}

\end{document}